The reference pipeline is run from the repository root:
\begingroup\small
\begin{verbatim}
python scripts/reproduce_report.py
\end{verbatim}
\endgroup
It executes both fitting modules with the supplied reference configurations,
CPU computation, one computational thread, and seed 7. The Neural SDE uses
1000 Monte Carlo paths and maximum normalized-time step 0.001. The logistic
SDE uses 5000 paths and maximum original-time step 0.05. Both use a central
90 percent process-noise interval. Each experiment saves its effective
arguments, package versions, input and source hashes, numerical settings, and
simulation diagnostics in \path{run_metadata.json} alongside its CSV outputs.
The transition objective remains the single-observation-step Gaussian
Euler-Maruyama approximation; forward integration refinement does not make
that likelihood exact.

The script writes the fitting and figure outputs to:
\begin{quote}
\small
\path{results/neural_sde_comparison_outputs}\\
\path{results/logistic_pinn_outputs}\\
\path{results/all_model_comparison_outputs}
\end{quote}
It generates
\begin{quote}
\small
\path{docs/generated/model_comparison_tables.tex}\\
\path{docs/generated/fit_values.tex}
\end{quote}
and a combined metrics CSV from the saved experiment results. The report
inputs the generated tables and values
directly. The script does not overwrite the existing compiled report PDF.
An isolated run can instead specify:
\begingroup\small
\begin{verbatim}
python scripts/reproduce_report.py \
  --output-root outputs/verification \
  --table-dir outputs/verification/tables
\end{verbatim}
\endgroup

To regenerate tables alone from the saved reference outputs:
\begingroup\small
\begin{verbatim}
python -m kefir_models.report_tables
\end{verbatim}
\endgroup
To rebuild the combined figure alone:
\begingroup\small
\begin{verbatim}
kefir-plot-all-models \
  --neural-comparison-csv \
    results/neural_sde_comparison_outputs/water_kefir_neural_dynamics_comparison.csv \
  --pinn-comparison-csv \
    results/logistic_pinn_outputs/water_kefir_logistic_pinn_comparison.csv \
  --pinn-dense-csv \
    results/logistic_pinn_outputs/water_kefir_logistic_pinn_dense.csv \
  --output-dir results/all_model_comparison_outputs \
  --interval-level 0.90
\end{verbatim}
\endgroup

Simulation uses a uniform internal grid determined by the endpoints and
maximum step. Requested values between grid points linearly interpolate
that approximation. Observation-grid and dense-grid logistic summaries use
the same random seed and integration grid. Thus requesting more plotting
points consumes no extra random draws and leaves shared-time predictions
unchanged. Positivity projection is a numerical modification of the logistic
Euler scheme; the number of projected values is saved and must be checked
under step refinement.

The recorded refinement and capacity-sensitivity checks can be repeated after
fitting with:
\begingroup\small
\begin{verbatim}
OMP_NUM_THREADS=1 MKL_NUM_THREADS=1 PYTHONPATH=src \
  python scripts/check_simulation_refinement.py --capacity-sensitivity
\end{verbatim}
\endgroup
Their CSV outputs are under \path{results/validation}. These finite checks
assess numerical sensitivity and Monte Carlo variability; they do not
establish asymptotic convergence or forecast calibration.

The five reference simulations hold fitted parameters fixed. The controlled
forecast experiment and Bayesian observation models are separate extensions:
\begingroup\small
\begin{verbatim}
pip install -e ".[bayes]"
python scripts/run_rev07.py
python scripts/run_bayesian_iteration.py
python scripts/run_structural_iteration.py
PYTHONPATH=src python scripts/build_iteration_tables.py
\end{verbatim}
\endgroup
Their outputs are saved under \path{results/rev07} and
\path{results/bayesian}, with structural alternatives under
\path{results/structural}. The structural runner reuses the upstream OU
posterior archives; run the Bayesian iteration first on a fresh checkout.
\texttt{--reuse-validated} verifies saved source hashes before reusing
the new structural chains. Tables and numerical summaries are generated from
those saved artifacts; the reference point-estimate table is retained as its
own comparison. Fixed seeds support repeatability within a recorded
environment; they do not promise bitwise equivalence across devices or
package versions. Compilation is performed separately from the experiment
pipeline so that an existing user PDF is preserved until explicitly rebuilt.

Readiness for independent-unit validation can be audited with:
\begingroup\small
\begin{verbatim}
PYTHONPATH=src python scripts/check_prospective_validation.py
PYTHONPATH=src python scripts/build_prospective_tables.py
PYTHONPATH=src python scripts/check_predictive_moments.py
\end{verbatim}
\endgroup
The draft specification is \path{configs/prospective_validation.json};
\path{data/metadata/reference.json} retains unconfirmed facts as null entries.
The audit and generated readiness CSVs are stored in
\path{results/prospective}. A readiness report is not a validation result.

The separate moment-controlled prior experiment is reproduced with:
\begingroup\small
\begin{verbatim}
python scripts/run_tempered_iteration.py
\end{verbatim}
\endgroup
It saves 40 refits under \path{results/tempered}, including diagnostics of
conditional response-mean observables. \texttt{--reuse-validated} verifies
the data and source hashes before reusing those new-prior chains. It does
not reuse old-prior samples as new-prior fits. The prior normalizers,
median calibration, forecast summaries, and generated LaTeX tables are
computed from the new case artifacts. Existing report PDFs are preserved.
